\subsection{What the Gaussian relaxation route requires}\label{round26:sec:ym-relaxation}
The internal-fibre framework supplied by James Montgomery Lockwood~\cite{Loc26Fibres} motivates the use of averaging projections and their spectra. The statements here test one such spectral route in an explicit finite-dimensional Gaussian model.
\subsection{The full Gaussian heat bath}


The conditional criterion of Corollary~\ref{cor:fibre-heatbath} asks for a
single-coordinate heat-bath gap $\lambda_b$ with
\[
 b\min(1,\lambda_b)/\log(2+\beta_b)\longrightarrow\infty.
\]
We show that the exact
gauge-fixed Gaussian fibre used in that manuscript violates this rate for
every positive sequence $\beta_b$. The Gaussian force covariance can still
have a vanishing absolute row sum. Thus estimates for the whole heat bath
and for the marked forces require different quantitative arguments.

\begin{lemma}\label{round26:linear}
Let $A$ be a positive definite symmetric matrix indexed by a finite set
$I$, let $D=\operatorname{diag}(A_{ii})$, and let $\nu$ be the centred
Gaussian law of covariance $(\beta A)^{-1}$, with $\beta>0$.
For $P_i=\mathbb E_\nu[\,\cdot\mid x_j,\ j\ne i]$, let
$L=\sum_i(I-P_i)$ and let $\lambda_{\rm hb}$ be its optimal variance gap.
For every nonzero vector $m\in\R^I$,
\[
 \lambda_{\rm hb}\le\frac{m^{\mathsf T}Am}{m^{\mathsf T}Dm}.
\]
\end{lemma}
\begin{proof}
Set $f(x)=(Dm)^{\mathsf T}x$. Gaussian conditioning gives conditional
variance $1/(\beta A_{ii})$ for the coordinate $x_i$. Therefore
\[
 \langle f,Lf\rangle_\nu=\frac{m^{\mathsf T}Dm}{\beta},\qquad
 \Var_\nu f=\frac{m^{\mathsf T}DA^{-1}Dm}{\beta}.
\]
The Cauchy--Schwarz inequality applied to $A^{1/2}m$ and $A^{-1/2}Dm$
gives
\[
 (m^{\mathsf T}Dm)^2\le(m^{\mathsf T}Am)(m^{\mathsf T}DA^{-1}Dm).
\]
Use $f$ in the variational definition of the gap.
\end{proof}

Here are the lattice definitions needed to apply the lemma.
Fix integers $b,L\ge3$, and let $\Lambda=(\Z/(bL)\Z)^4$. A link is
$(x,\mu)$ from $x$ to $x+e_\mu$, with $1\le\mu\le4$.
For $y\in(\Z/L\Z)^4$ and each $\mu$, include in the path set $P$ all
the $b$ links from $by$ to $b(y+e_\mu)$ in direction $\mu$. Write $E_0$
for the links outside $P$, and $\Sigma$ for all sites on these paths.
For $m\in\R^{E_0}$, let $\widetilde m$ be its extension by zero to $P$.
Let $d_1$ take a link field to its oriented sums around the plaquettes,
and let $d_0$ be the gradient on site fields. Define
\[
 K_0=(d_1^{\mathsf T}d_1)_{E_0,E_0},\qquad
 B\phi=(d_0\phi)|_{E_0},\qquad A=K_0+BB^{\mathsf T},
\]
where the domain of $B$ consists of site fields vanishing on $\Sigma$.
The manuscript uses this same $A$ to compute its Gaussian fibre kernel.

\begin{theorem}\label{round26:slow}
The matrix $A$ is positive definite, and its diagonal entries lie in
$[6,8]$. For the Gaussian probability measure of precision $\beta A$,
with the generator normalized as the sum of the single-link heat-bath
generators, its gap satisfies
\[
 0<\lambda_{b,L,\beta}\le\frac{1}{b^3-1}
 \qquad(b,L\ge3,\ \beta>0).
\]
The upper bound is uniform in the coarse torus size and the coupling.
The test vector producing it belongs to the orthogonal complement of
the gauge gradients $\operatorname{ran}B$.
\end{theorem}
\begin{proof}
First we justify positivity. If $d_1\widetilde m=0$, lift the periodic
closed form $\widetilde m$ to $\Z^4$. Its integral along paths defines a
potential $\phi$ with $d_0\phi=\widetilde m$, because every closed lattice
loop in $\Z^4$ is generated by plaquette boundaries. The lifted skeleton
is connected and $\widetilde m$ vanishes on its links. Choose the additive
constant so that $\phi=0$ there. Its difference under translation by a
period has zero gradient and vanishes on the skeleton, so $\phi$ is
periodic. Thus $m=B\phi$ and $\ker K_0=\operatorname{ran}B$.
If $m^{\mathsf T}Am=0$, both $d_1\widetilde m=0$ and $B^{\mathsf T}m=0$.
Then $m=B\phi$ and $\|m\|^2=\langle\phi,B^{\mathsf T}m\rangle=0$.
This proves positive definiteness. Every fine link is contained in six
plaquettes, giving $(K_0)_{ii}=6$. The entry $(BB^{\mathsf T})_{ii}$ is
the number of endpoints of that link outside $\Sigma$, which is $0$, $1$
or $2$. This proves the diagonal bounds.

Now take $m$ to be $1$ on free links of direction $1$ and zero on the
other free links. There are $(bL)^4$ links of direction $1$, and exactly
$bL^4$ of these belong to $P$. Consequently
\begin{equation}\label{round26:norm}
 \|m\|^2=(b^4-b)L^4.
\end{equation}
At a site outside $\Sigma$, every incident link is free, because every
path link has both endpoints in $\Sigma$. The incoming and outgoing
direction-$1$ values of $m$ both equal $1$. Its divergence at every such
site is therefore zero. This proves $B^{\mathsf T}m=0$.

The field equal to $1$ on all direction-$1$ links and zero in the other
directions has zero curl. The extension $\widetilde m$ differs from it
only on the $bL^4$ direction-$1$ links in $P$. No plaquette contains two
of those deleted links. For opposite direction-$1$ sides of a plaquette,
one transverse coordinate differs by $1$, whereas a direction-$1$ path
requires every transverse coordinate to be divisible by $b$. Each deleted
link belongs to exactly six plaquettes. Thus exactly $6bL^4$ plaquettes
have curl $1$ or $-1$, and all others have zero curl. We obtain the exact
energy identity
\begin{equation}\label{round26:energy}
 m^{\mathsf T}Am=\|d_1\widetilde m\|^2+
                 \|B^{\mathsf T}m\|^2=6bL^4.
\end{equation}
Since $D\ge6I$, Lemma~\ref{round26:linear} and \eqref{round26:norm}--\eqref{round26:energy} give
\[
 \lambda_{b,L,\beta}\le\frac{6bL^4}{6(b^4-b)L^4}
       =\frac1{b^3-1}.
\]
The positive lower bound at each fixed finite volume follows from the
Gaussian projection argument given below. Only the displayed upper bound
is used in the conclusion about the required rate.
\end{proof}

For completeness, the positive lower bound in Theorem~\ref{round26:slow} has an
elementary conditional-variance proof that avoids derivative estimates.
Let $Y$ be a standard Gaussian vector and write $x=(\beta A)^{-1/2}Y$.
The conditioning subspace for coordinate $i$ has one-dimensional
orthogonal complement spanned by a unit vector $u_i$. The vectors $u_i$
span the first Gaussian chaos. Put $M=\sum_i u_i u_i^{\mathsf T}$ and
$\eta=\lambda_{\min}(M)>0$. On the symmetric tensor power of order $q\ge1$,
the Gaussian conditional projection is $(I-u_i u_i^{\mathsf T})^{\otimes q}$.
For any orthogonal projection $Q$, the tensor factors commute and
\[
 I-Q^{\otimes q}\ge\frac1q\sum_{j=1}^q (I-Q)_j.
\]
It follows that
\[
 \sum_i\left[I-(I-u_i u_i^{\mathsf T})^{\otimes q}\right]\ge\eta I.
\]
The orthogonal decomposition into Gaussian chaoses therefore gives
\[
 \langle f,Lf\rangle\ge\eta\Var(f)\qquad(f\in L^2).
\]

This example also preserves the marked-force covariance from the quotient
Gaussian fibre. For a mark $a\in P$ and its force coefficient
$\psi_a=(d_1^{\mathsf T}d_1\delta_a)|_{E_0}$ one has
$\psi_a\perp\operatorname{ran}B$, since $d_1d_0=0$. On that orthogonal
complement, $A^{-1}=K_0^+$, where $K_0^+$ is the inverse of $K_0$ on its
kernel complement. Hence the force covariance computed using $A^{-1}$ is
exactly the one computed by the manuscript's quotient Gaussian measure.
Its previously established row-sum estimate is $O(\beta_b(1+\log b)/b)$.
For example, this tends to zero when $\beta_b=O(\log b)$.

In contrast, Theorem~\ref{round26:slow} implies
\[
 \frac{b\min(1,\lambda_{b,L,\beta_b})}{\log(2+\beta_b)}
 \le\frac{b}{(b^3-1)\log2}\longrightarrow0.
\]
Thus the required whole-fibre relaxation rate fails even in this quadratic
model, where the force covariance is known and small along logarithmic
couplings. The comparison with the compact Wilson model still requires
estimates of its own. The calculation identifies a quantitative
obstruction to deriving the sufficient full heat-bath rate from Gaussian
force decorrelation alone.

\subsection{The spectral response of marked forces}


Fix integers $b,L\ge3$, let $\Lambda=(\Z/(bL)\Z)^4$, and let
$N=L^4$. A directed link $(x,\mu)$ belongs to the straight-path skeleton
$P$ when $x_\rho$ is divisible by $b$ for every $\rho\ne\mu$.
The sites $\Sigma$ on the skeleton have at most one coordinate that is
nonzero modulo $b$. Put $E_0=E\setminus P$. A vector on $E_0$ is
extended by zero on $P$ when it occurs in a plaquette sum.
Let $d_0$ be the site-to-link gradient and $d_1$ the link-to-plaquette
curl. Define
\[
 \begin{gathered}
 K=d_1^{\mathsf T}d_1,\qquad K_0=K_{E_0,E_0},\qquad
 B\phi=(d_0\phi)|_{E_0},\\
 A=K_0+BB^{\mathsf T},\qquad
 D=\operatorname{diag}(A_{ii}),\qquad H=D^{-1/2}AD^{-1/2}.
 \end{gathered}
\]
Here $\phi$ is a site field off $\Sigma$, extended by zero on $\Sigma$.
The kernel of $K_0$ is $\operatorname{ran}B$. Indeed, a closed zero-on-$P$
link field lifts to a gradient on $\Z^4$; its potential is constant on the
connected lifted skeleton and can be set to zero there. Its period
differences are then zero as well. Consequently $A$ is positive definite
and $6I\le D\le8I$, since a free link lies in six plaquettes and has at
most two endpoints off $\Sigma$.

The mark of the coarse path $(y,\mu)$ is
$a=(by+\lfloor b/2\rfloor e_\mu,\mu)$. For a mark put
\[
 \psi_a=(K\delta_a)|_{E_0},\qquad h_a=D^{-1/2}\psi_a.
\]
On the Gaussian law $\nu$ of covariance $(\beta A)^{-1}$, $\beta>0$, the
centred marked force is $F_a(x)=\beta\psi_a^{\mathsf T}x$.
The whole-fibre generator is $\mathscr L=\sum_{i\in E_0}(I-P_i)$,
where $P_i$ is conditional expectation given all coordinates except $i$.
Write $\mathscr P_t=\exp(-t\mathscr L)$.
Let $\mathcal M$ denote the $N$ marks of direction $1$, and put
\[
 p=\sum_{a\in\mathcal M}\psi_a,\qquad h=D^{-1/2}p,
 \qquad F=\sum_{a\in\mathcal M}F_a.
\]

\begin{theorem}
For the spectral measure $\mu_h$ of $h$ under the positive matrix $H$,
\[
 \int s^{-1}\,d\mu_h(s)=6N,\qquad
 \int s^{-2}\,d\mu_h(s)\ge\frac{6N(b^3-1)}b.
\]
In particular,
\begin{equation}\label{round26:iat}
 \frac{\displaystyle\int_0^\infty\Cov_\nu(F,\mathscr P_tF)\,dt}
      {\Var_\nu F}\ge\frac{b^3-1}b.
\end{equation}
Every local mark $a\in\mathcal M$ has a nonzero spectral component under
$H$ at some rate at most $b/(b^3-1)$. More quantitatively, for
$\theta=2b/(b^3-1)$ and its spectral projector $\Pi_\theta$, each satisfies
\begin{equation}\label{round26:overlap}
 \|\Pi_\theta H^{-1}h_a\|^2\ge\frac{3(b^3-1)}{bN}.
\end{equation}
The normalized integrated autocorrelation of each local force is at least
$(b^3-1)/(bN)$.
\end{theorem}

\begin{proof}
For any source $q$, Gaussian conditioning on coordinate $i$ gives
\[
 (I-P_i)(q^{\mathsf T}x)=\frac{q_i}{A_{ii}}(Ax)_i.
\]
The coefficient of a linear observable therefore evolves by
$\exp(-tAD^{-1})$. Put $u=D^{-1/2}q$. Then
\begin{equation}\label{round26:spectral}
 \Cov_\nu(\beta q^{\mathsf T}x,
             \mathscr P_t(\beta q^{\mathsf T}x))
   =\beta u^{\mathsf T}H^{-1}e^{-tH}u.
\end{equation}
In particular this correlation is nonnegative. Its value at zero is
$\beta\int s^{-1}d\mu_u$, and its time integral is
$\beta\int s^{-2}d\mu_u$.

For the source $p$, let $w=A^{-1}p$. Each $\psi_a$ is orthogonal to
$\operatorname{ran}B$, since $d_1d_0=0$. Hence $w$ is orthogonal to
$\operatorname{ran}B$ and $K_0w=p$.
Prescribe values $1$ on all marks in $\mathcal M$, and $0$ on the other
path links. Call that path field $v$. Different middle marks are in no
common plaquette. The full Maxwell energy, with free field $m$, is
therefore
\[
 \|d_1(v+m)\|^2=6N+2m^{\mathsf T}p+m^{\mathsf T}K_0m.
\]
Its minimum occurs at $m=-w$ modulo the gauge kernel.
The minimum is zero. To see this, start with the closed field equal to
$1/b$ on every direction-$1$ link and zero otherwise. Along each
direction-$1$ skeleton path, choose a potential that is zero at both
ends, with increments $-1/b$ except for the marked increment $1-1/b$.
Set the potential equal to zero on the other paths, which is consistent
at their common coarse endpoints, and extend it to the remaining sites.
Adding its gradient gives the required values on $P$ and preserves zero
curl. It follows that
\begin{equation}\label{round26:first}
 p^{\mathsf T}A^{-1}p=p^{\mathsf T}w=6N.
\end{equation}

The minimizing full field $v-w$ is closed. Its integral around the fine
direction-$1$ cycle with transverse coordinates zero is $L$, because
that cycle consists of $L$ skeleton paths, with a single unit mark on
each. Closedness makes its integral equal to $L$ on every parallel fine
cycle, by summing the plaquettes between two adjacent cycles.
There are $(bL)^3-L^3=L^3(b^3-1)$ cycles whose transverse coordinates
are not all divisible by $b$. Every link on such a cycle is free. Each
cycle has $bL$ links, so Cauchy--Schwarz yields a squared free-field norm
of at least $L^2/(bL)=L/b$ on each of these disjoint cycles. Therefore
\begin{equation}\label{round26:second}
 \|w\|^2\ge\frac{N(b^3-1)}b,
 \qquad w^{\mathsf T}Dw\ge\frac{6N(b^3-1)}b.
\end{equation}
Since $H^{-1}h=D^{1/2}w$, equations \eqref{round26:first} and \eqref{round26:second}
are the stated spectral moment identities. Equation \eqref{round26:spectral}
then gives \eqref{round26:iat}.

The ratio of the first inverse moment to the second is the average of
$s$ against the probability measure proportional to $s^{-2}d\mu_h(s)$.
It is at most $b/(b^3-1)$. Thus $h$ has a nonzero spectral component at
or below that rate. Translations by $b$ in each coordinate preserve $A$
and $D$, and act transitively on $\mathcal M$. If one $h_a$ lacked every
component at or below this rate, they all would, and so would their sum
$h$. This proves the local support assertion.

For the quantitative projector bound, the portion of the second inverse
moment above $\theta$ is at most $\theta^{-1}\int s^{-1}d\mu_h$.
Using \eqref{round26:first} and \eqref{round26:second}, this gives
\[
 \|\Pi_\theta H^{-1}h\|^2\ge\frac{3N(b^3-1)}b.
\]
On the other hand, $h=\sum_a h_a$ and the squared norm of a sum of $N$
vectors is at most $N$ times the sum of their squared norms.
Translation invariance makes the $N$ local norms equal, proving
\eqref{round26:overlap}.
The same argument without the projector shows
\[
 \int_0^\infty\Cov_\nu(F_a,\mathscr P_tF_a)\,dt
    \ge\frac{6\beta(b^3-1)}{bN}.
\]
Finally $\Var_\nu F_a\le6\beta$. Indeed, minimizing the nonnegative
Maxwell quadratic form with just that mark set equal to $1$ gives
$6-\psi_a^{\mathsf T}A^{-1}\psi_a\ge0$. Divide by the local variance to
obtain the last claim.
\end{proof}

For fixed coarse side $L$, the local normalized time integral grows at
least like $b^2$. Moreover the support assertion holds for every $L$.
Consequently the spectral subspace generated by all local marked-force
observables under the exact heat-bath semigroup cannot have a gap with
\[
 b\lambda_b/\log(2+\beta_b)\longrightarrow\infty
\]
in this Gaussian model. The earlier whole-fibre slow-mode bound is
$\lambda_b\le1/(b^3-1)$. The present argument does not assume that the
forces are orthogonal to every slow mode, and proves directly that such
orthogonality fails.

These temporal estimates coexist with the manuscript's Gaussian
absolute force row-sum estimate $O(\beta_b(1+\log b)/b)$. A Wilson
comparison therefore needs control of the spatially summed force
covariances, or a semigroup response estimate retaining both spatial
decay and the sizes of spectral weights. A uniform spectral
support gap for the force sector alone cannot supply the required rate
through this quadratic approximation. A corresponding estimate for the
compact nonabelian Wilson fibre remains to be proved.

\begin{figure}[tb]
\centering
\begin{tikzpicture}[>=Latex,font=\scriptsize,
 box/.style={draw,rounded corners,align=center,text width=39mm,minimum height=10mm,fill=blue!5}]
\node[box] (g) at (0,0) {Gaussian constrained fibre};
\node[box] (h) at (0,-1.3) {full heat-bath gap\\$\lambda_b\le(b^3-1)^{-1}$};
\node[box] (f) at (0,-2.7) {marked-force response\\at some $s\le b/(b^3-1)$};
\node[box,fill=green!5] (r) at (0,-4.1) {spatial covariance row sum\\$O(\beta\log b/b)$};
\draw[->] (g)--(h); \draw[->] (h)--(f);
\draw[->] (g.east)--(2.6,0)--(2.6,-4.1)--(r.east);
\end{tikzpicture}
\caption{The Gaussian comparison has small spatial covariance row sums while both full and force-generated relaxation can be slow. A spectral-gap argument alone does not supply the Wilson covariance estimate.}
\label{round26:fig:ym-relaxation}
\end{figure}
